% ============================================================================
% 元数据配置 - config.tex
% 包含论文标题、作者信息、摘要、关键词等元数据
% ============================================================================

% --- 论文标题 ---
\newcommand{\papertitle}{基于风险时空异质性的城市低空无人机航路规划}

% --- 作者信息 ---
\newcommand{\authorname}{xxx, xxx, xxx}
\newcommand{\authoremail}{xxx@example.com}
\newcommand{\authoraddress}{xxxx大学xxx院，城市邮编}
\newcommand{\correspondingnote}{通讯作者}

% --- 摘要内容 ---
\newcommand{\abstracttext}{
城市低空空域的无人机运行在提升公共与商业服务效能的同时，不可避免地对地面第三方构成了安全与社会滋扰风险，这也成为制约其大规模应用的关键因素。
现有风险感知路径规划研究多假设理想静态环境与经验恒定坠机概率，未能有效刻画城市复杂微气象、建筑阻隔以及风险时空特征所导致的动态风险变化。
因此，本文提出一种多维时空耦合的第三方风险（TPR）综合成本评估模型。该模型突破常数假设，结合比例风险模型引入风、雨及城市峡谷因子重构动态坠机概率；
利用时空潮汐引力模型模拟动态人口与车辆暴露密度；并从物理致死、建筑财产损失及社会时空噪声敏感度等多维度实现了风险代价的更为精确的量化。
针对高维动态风险模型导致的维度灾难与搜索空间爆炸，本文结合分布估计算法（EDA）、K-means 空间聚类与 CostA 算法，构建了一种分层混合路径寻优策略。利用 EDA 进行全局探索，通过 K-means 提取动态低风险走廊实现空间降维，最终依靠 CostA* 完成精确的局部轨迹规划。
模拟实验和真实城市场景实验表明，所提方法能在极小的绕行代价下，显著降低核心城区的物理撞击风险与敏感时段噪声滋扰。较之传统静态启发式规划，该框架不仅具备更高的复杂环境适应性，且能根据城市时间节律与多目标权重偏好，生成高可解释性的安全飞行路径。
}

% --- 关键词 ---
\newcommand{\paperkeywords}{
第三方风险 \sep 时空异质性 \sep 无人机 \sep 城市低空 \sep 路径规划
}

% --- MSC 分类代码 ---
\newcommand{\msccodes}{90Cxx \sep 68Txx}